\scanpage{4}
\mspage{-3.2-}

\margin{\handwritten{\unsure{11|''}}}

\[
\handwritten{\varpi}\colon\mathrm{SL}(n+1)\handwritten{\longrightarrow}\mathrm{PGL}(n)
\]

is the canonical isogeny, then with respect to the induced action of $\mathrm{SL}(n+1)$, $\mathrm{o}_{\handwritten{\mathbb P}}(1)$ admits an $\mathrm{SL}(n+1)$-linearization. This follows because $\mathrm{SL}(n+1)$ acts on the affine cone over $\handwritten{\mathbb P}_n$, compatibly with the given action of $\mathrm{PGL}(n)$ on $\handwritten{\mathbb P}_n$. This linearization induces an $\mathrm{SL}(n+1)$-linearization of $L_i$ for each $i$ (for the induced action of $\mathrm{SL}(n+1)$ ), and therefore such a linearization for each product

\[
\handwritten{\bigotimes}_{\handwritten{i=0}}^{\handwritten{m}} L_i^{\handwritten{\otimes r_i}}.
\]

On the other hand, the $(n+1)^{\mathrm{st}}$ power of the $\mathrm{SL}(n+1)$-linearization of $\underline{\mathrm{o}}_{\handwritten{\mathbb P}}(1)$, which is a linearization of $\underline{\mathrm{o}}_{\handwritten{\mathbb P}}(n+1)$, is induced, via $\handwritten{\varpi}$, from a $\mathrm{PGL}(n)$-linearization (cf. Chapter 1, \S2). Therefore, the $\mathrm{SL}(n+1)$-linearization of

\[
\handwritten{\bigotimes}_{\handwritten{i=0}}^{\handwritten{m}} L_i^{r_i}
\]

is a $\mathrm{PGL}(n)$-linearization if $n+1$ divides each $r_i$.

In any case, by means of these linearizations we can define invariant sections of all the sheaves $\handwritten{\bigotimes}_{\handwritten{i=0}}^{\handwritten{m}} L_i^{r_i}$. To construct such invariant sections, let $X_{\unsure{o|0}},\ldots,X_n$ be the canonical sections of $\underline{\mathrm{o}}_{\handwritten{\mathbb P}}(1)$ on $\handwritten{\mathbb P}_n$. Let

\[
X_i^{(j)}=p_j^*(X_i)
\]

be the induced sections of $L_j$.

\uline{Definition} \uline{3.1}: For all sequences $\handwritten{\alpha}_o,\ldots\handwritten{\alpha}_n$ of integers such that

% archive sheet number: 3
